% !TeX program = xelatex
% UTF-8; standalone source. Compile with XeLaTeX twice.
\documentclass[UTF8,a4paper,zihao=-4,fontset=fandol]{ctexart}
\usepackage[margin=2.5cm]{geometry}
\usepackage{fancyhdr,xurl}
\usepackage[hidelinks]{hyperref}
\pagestyle{fancy}
\fancyhf{}
\fancyhead[L]{\small 科技创新与青年担当}
\fancyhead[R]{\small 程泽凯\quad 2023112621}
\fancyfoot[C]{\small\thepage}
\setlength{\headheight}{16pt}
\setlength{\parindent}{2em}
\setlength{\parskip}{0pt}
\linespread{1}
\ctexset{section={format=\large\bfseries,beforeskip=14pt,afterskip=7pt}}
\newcommand{\bodyfont}{\fontsize{12bp}{18bp}\selectfont}
\begin{document}
\bodyfont
\begin{center}
  {\zihao{2}\bfseries 把青春写进创新链}\par\vspace{5pt}
  {\large 从专业学习走向科技报国}\par\vspace{12pt}
  姓名：程泽凯\qquad 学号：2023112621
\end{center}
\vspace{6pt}
\noindent\textbf{摘要：}中国共产党成立一百零五周年，为当代大学生理解个人成长与国家发展的关系提供了历史坐标。青年参与强国建设，既要明确价值方向，也要建立能够解决真实问题的知识结构和实践能力。本文以科技创新为切入点，讨论地方发展需求、校园学习资源与专业训练之间的联系，主张通过问题发现、基础学习、协作验证和成果交付，把报国志向转化为可持续的行动。同时，技术实践必须坚持求真、尊重规则，关注应用对象的实际需要，使创新既有技术质量，也有社会价值。

\noindent\textbf{关键词：}青年担当；科技创新；专业学习；实践能力

% BODY_START
\section{历史坐标中的青年选择}
2026年是中国共产党成立105周年。从嘉兴南湖上的一叶红船出发，中国共产党团结带领人民走过艰难而壮阔的奋斗历程。理解这一历史，不能只记住若干时间节点，还应看到贯穿其中的责任意识：在国家和民族面临重大问题时，总有人把个人前途与共同命运联系起来，把对未来的期待变成持续的组织和行动。对于当代大学生，这种联系依然存在，只是担当的具体内容有了新的时代条件。

在革命年代，青年面对的是民族独立和人民解放的紧迫任务；在今天，建设社会主义现代化强国要求我们处理发展质量、科技能力、公共服务和社会公平等问题。青年不必机械复制前人的生活方式，却应继承他们面对现实、主动承担的精神。科技进步给大学生提供了更广阔的学习渠道，也提出了更高要求。会使用某种工具，并不自动意味着具备创造能力；能够表达爱国情感，也不自动意味着已经完成报国责任。

我理解的青春担当，是愿意对一个有价值的问题负责，并为解决它付出长期努力。一个学生的影响范围可能有限，但他能否认真完成学习任务、如实记录实验结果、主动纠正设计缺陷，关系到未来是否能承担更复杂的工作。科技报国并非毕业以后才开始的宏大事业，它首先体现为今天怎样学习、怎样判断、怎样与他人合作。这种从日常行为出发的理解，使国家发展的目标能够进入个人成长的实际过程。

\section{让国家需要成为发现问题的尺度}
科技创新为什么与强国建设紧密相关？因为现代社会的生产、治理和生活越来越依赖稳定而可靠的技术系统。关键环节缺乏自主能力，发展就容易受制于外部条件；技术成果不能进入实际场景，研究投入也难以形成充分的社会收益。2024年全国科技大会相关讲话强调科技创新与国家需求的联系，并讨论教育、科技、人才的相互支撑关系\cite{science}。对大学生而言，这提供了一种理解专业学习的视角：知识价值既在于解释世界，也在于帮助社会解决问题。

国家需要并不是一个脱离地方生活的抽象概念。它可以体现为企业希望减少生产故障，乡村希望降低农业经营成本，社区希望改善办事体验，公共机构希望提高信息服务的可靠性。不同省份的资源禀赋、产业结构和发展阶段并不相同，同一个技术方案也可能面对不同约束。因此，关注省情首先应关注真实需求，理解当地哪些问题突出、哪些条件已经具备、哪些限制短期内难以改变，而不能只根据热门词语选择研究方向。

这种问题意识可以通过有组织的观察培养起来。学生在查阅地方公开资料、参加实践或开展访谈时，应先区分事实、解释和主张。企业说某个环节效率低，究竟是技术性能不足，还是流程安排不合理？使用者觉得系统难用，是界面复杂，还是培训不足？如果问题判断错误，即使方案本身很先进，也可能没有效果。青年把专业所学与地方发展连接起来，第一步通常是认真提问，而不是急于展示一个已经想好的答案。

选择课题时还应考虑自己的能力边界。国家发展的任务很大，学生项目的时间、资源和经验却有限。比较稳妥的方式，是从大问题中切出一个可以验证的小问题。例如讨论公共信息服务时，可以先研究一项操作是否清楚、一个数据字段是否准确、一种提醒是否容易理解。这样做并不削弱课题的意义，反而可以让责任落实到具体对象，使学习成果具有清晰的评价标准。

\section{以校园学习为创新能力的起点}
学校是大学生走向社会的准备空间。结合校情，并不意味着简单罗列学校的荣誉和口号，而是要弄清自己能够利用哪些课程、图书资料、实验条件和指导资源，同时认识这些资源的限制。课程中的理论、实验中的操作、讨论中的反馈，本来就是一个相互支持的学习系统。学生若只把它们视为获取学分的独立任务，就容易错过知识之间的联系，也难以形成完整的问题解决能力。

首先应建立扎实的基础。不同专业所需的基础知识有所区别，但严谨推理、资料阅读、清楚表达和规范操作具有共同意义。学习一个概念时，既要理解它能解释什么，也要理解它在什么条件下不适用。遇到公式或模型，不应只记住最终形式，还应追问假设从哪里来、变量代表什么、改变条件会产生什么影响。基础学习的收益不一定立即显现，却决定了面对新问题时能否作出可靠判断。

其次应让课程作业形成连续积累。一门课结束之后，代码、笔记、调研材料和设计记录不应立即失去用途。可以把已经学会的方法整理成个人知识档案，标明适用场景、常见错误和尚未理解的部分。下一次遇到相近问题时，先检查旧方法是否适用，再考虑新增工具。这样的学习方式能够减少对零散技巧的依赖，也能帮助学生看见自己的进步究竟发生在哪些方面。

再次应重视跨专业沟通。现实问题往往同时涉及技术、成本、制度和使用习惯，一种专业视角很难覆盖全部条件。参与合作时，专业能力强的同学需要把判断依据讲清楚，熟悉社会情境的同学也需要说明需求来自何处。有效协作的前提不是大家使用同样的术语，而是能够围绕同一个问题核对信息、讨论取舍。对大学生而言，学会解释自己的方法，与学会使用方法同样重要。

最后应把独立思考与合理使用工具结合起来。智能工具可以帮助检索、整理和生成初步方案，但不能替代学生对事实和逻辑的核查。面对自动生成的结论，应检查来源是否可靠、推导是否成立、内容是否适合当前问题。若一个学生只能展示工具给出的结果，却无法说明为什么选择它，也无法解释失败原因，那么得到的更多是一份临时答案，而不是可以迁移的能力。

\section{把实干落实到完整的实践过程}
从专业学习走向科技实践，关键在于建立完整过程，而不是停留在概念展示。一个项目至少需要回答四个问题：服务谁、解决什么、如何证明有效、由谁继续维护。只有把这些问题说清楚，才能判断技术选择是否合理。对于学生团队，项目规模可以小，但过程应当完整。明确目标、形成计划、记录变化、接受反馈，都属于创新能力的一部分，不应被视为正式工作之外的附属事项。

在准备阶段，应先了解现有办法，而不是预设一切都需要重新发明。阅读公开资料、学习已有方案、访谈潜在使用者，可以帮助确定真正值得改进的环节。如果现有方法已经能够低成本完成任务，就需要解释新方案增加了什么价值。创新可以是某种性能改善，也可以是流程更清楚、操作更方便或维护更容易。把创新等同于使用最新工具，容易忽略解决问题本身。

在验证阶段，应设置合理的对照与评价条件。对于能够量化的问题，可以比较处理时间、错误情况或资源消耗；对于不适合简单量化的问题，可以结合使用反馈和规范要求进行分析。无论采用什么方式，都应说明条件是否一致、材料是否足够、结果能否重复。一个偶然成功的展示不能替代稳定效果，某个指标变好也不代表整个方案更好。诚实说明限制，是让成果值得信任的必要步骤。

在协作阶段，应建立清晰的责任分工。谁负责收集资料，谁负责实现，谁负责核查，谁负责与使用者沟通，应在开始时有所约定。项目变化后也应及时调整，而不能让责任模糊地落在最积极的人身上。团队成员应尊重彼此的劳动，记录引用来源和贡献情况。遇到分歧时，应回到共同目标和事实依据上讨论，避免把方案争论变成对个人能力的简单评价。

在交付阶段，应把说明文件和后续支持纳入成果。一个作品即使能运行，如果别人不知道如何使用、无法处理异常或无法继续更新，其实际价值仍然有限。大学生可以练习写简明操作说明、整理必要材料、记录已知问题，并在条件允许时安排回访。把成果交到别人手里以后仍然关心它的使用情况，体现的是对服务对象的责任，也使短期课程项目有机会转化为持续改进的基础。

\section{以求真和责任守住创新的质量}
青年参与科技创新，还必须面对价值判断。技术是否先进与应用是否合理，是相互联系却不能混为一谈的两个问题。收集数据前应考虑是否确有必要，涉及个人信息时应尊重知情和选择，展示结果时应避免夸大能力。如果为了让项目更有吸引力而省略限制、包装数据，短期或许显得漂亮，长期却会损害使用者的利益和团队的信誉。求真不是创新速度的障碍，而是有效创新的基础。

面对困难，也应保持合理的奋斗观。实干不等于无限延长工作时间，更不是用忙碌掩盖方向不清。重要的是能否保持学习节奏、主动反馈问题、在失败后调整方法。大学生要兼顾身体健康和心理状态，承认能力形成需要过程。一次项目没有达到目标，可以成为认识条件和修正判断的机会；只有拒绝复盘、回避事实，才会让失败无法转化为经验。

落实到个人行动，可以以一个学期为单位建立计划：围绕课程夯实一项基础能力，围绕现实需求理解一个具体问题，围绕合作完成一次有记录的实践，再通过总结找出下一阶段的不足。评价这些行动时，不只看获奖或展示效果，还应看是否学会解释方法、是否能独立处理错误、是否给他人带来实际帮助。持续完成这样的循环，才能逐步形成参与更大任务的能力。

《新时代的中国青年》白皮书呈现了青年在不同领域参与社会发展、承担时代责任的多种方式\cite{youth}。这种多样性提醒我们，科技报国没有唯一的个人路线。有人适合基础研究，有人擅长工程实现，也有人能够把技术与公共服务连接起来。对当代大学生而言，重要的是找到能力与需要相交的位置，以诚实学习积累本领，以实际行动承担责任，让青春在强国建设和民族复兴的征程中形成坚实的贡献。
% BODY_END

\begin{thebibliography}{9}
\fontsize{10.5bp}{15.75bp}\selectfont
\bibitem{science} 习近平. 在全国科技大会、国家科学技术奖励大会、两院院士大会上的讲话[EB/OL]. 2024-06-24. \url{https://www.gov.cn/gongbao/2024/issue_11466/202407/content_6963180.html}.
\bibitem{youth} 国务院新闻办公室. 新时代的中国青年[EB/OL]. 2022-04-21. \url{https://www.mee.gov.cn/zcwj/gwywj/202204/t20220421_975555.shtml}.
\end{thebibliography}
\end{document}
